%% kapa: exam=1 round=23 year=2012 period=1 form=A subject=economics subject_ko=경제원론 seq=5 number=45
%% source: Q-Net 감정평가사 기출문제 (cst00309) artlSeq=5052578 file=raw/1008/5052578/2012년 제23회 감정평가사 제1차시험 기출문제(1교시 A형).pdf
%% conversion: kordoc 4.24.51 -> tools/gen1.py (auto); review=proofread-sonnet; flags=none
\begin{question}{45}
고정비용이 존재하고 노동만이 가변요소인 기업의 단기비용에 관한 설명으로 옳지 않은 것은?
\begin{choices}
  \item 단기평균고정비용 곡선은 언제나 우하향한다.
  \item 단기총평균비용은 단기평균가변비용과 단기평균고정비용의 합이다.
  \item 노동의 한계생산이 체감하면 단기한계비용 곡선은 우상향한다.
  \item 노동의 한계생산이 불변이면 단기총평균비용 곡선은 수평이다.
  \item 단기한계비용이 단기총평균비용보다 큰 경우 단기총평균비용은 증가한다.
\end{choices}
\end{question}
\answer{4}
